\section{To‘g‘ri chiziqli tekis harakat}

\begin{enumerate}

\item Qanday harakat to‘g‘ri chiziqli tekis harakat bo‘ladi?
\begin{enumerate}[label=\Alph*)]
\item Teng vaqt oraliqlarida bir xil masofaga ko‘chuvchi va trayektoriyasi to‘g‘ri chiziqdan iborat bo‘lgan jism harakati
\item Trayektoriyasi to‘g‘ri chiziqdan iborat harakat
\item Hech qanday tezlanishga ega bo‘lmagan harakat
\item Teng vaqt oraliqlarida teng masofalar o‘tuvchi jism harakati
\end{enumerate}

\item To‘g‘ri chiziqli tekis harakatda tezlik vektorining ...
\begin{enumerate}[label=\Alph*)]

\item moduli tekis oshib, yo‘nalishi o‘zgarmaydi
\item moduli va yo‘nalishi o‘zgarmas saqlanadi
\item moduli va yo‘nalishi o‘zgarib turadi
\item moduli o‘zgarmas bo‘lib, yo‘nalishi uzluksiz o‘zgarib turadi
\end{enumerate}

\item 918 km/h tezlik necha m/s bo‘ladi?
\begin{enumerate}[label=\Alph*)]
\item 260
\item 240
\item 250
\item 255
\end{enumerate}

\item Uzunligi 9 m bo‘lgan ingichka ip boshidan jo‘nagan chumoli ip oxiriga kelguncha 5 minut o‘tdi. 
Chumolining tezligini toping (cm/s).
\begin{enumerate}[label=\Alph*)]
\item 1,8
\item 3
\item 3,5
\item 4,5
\end{enumerate}

\item Uzunligi 200 m bo‘lgan poyezd 10 m/s tezlik bilan tekis harakatlanib, uzunligi 300 m bo‘lgan tunnelga kirib bormoqda. 
Necha sekunddan keyin poyezd tunneldan butunlay chiqib ketadi?
\begin{enumerate}[label=\Alph*)]
\item 50
\item 30
\item 25
\item 20
\item 10
\end{enumerate}

\item Birinchi avtomobil $4S$ masofani $2t$ vaqtda $2v$ tezlik bilan bosib o‘tdi. 
Ikkinchi avtomobil $3S$ yo‘lni $2t$ vaqtda bosib o‘tsa, uning tezligini toping.
\begin{enumerate}[label=\Alph*)]
\item $2.25v$
\item $2.5v$
\item $1.5v$
\item $0.5v$
\item $v$
\end{enumerate}

\item Magnitofon o‘ragichi $4\ \text{m/s}$ tezlik bilan $50\ \text{s}$ da tasmani o‘rab bo‘ldi. 
Agar o‘ragichning boshlang‘ich radiusi $3$ cm, oxirgi radiusi $7$ cm bo‘lsa, tasmaning qalinligini toping (mm). 
$\pi=3$ deb oling.
\begin{enumerate}[label=\Alph*)]
\item 0.5
\item 0.06
\item 0.1
\item 0.09
\item 0.08
\end{enumerate}

\item Uzunligi 300 m bo‘lgan, 36 km/h tezlik bilan harakatlanayotgan mashinalar qatori 60 m uzunlikdagi ko‘prikdan necha sekundda o‘tib ketadi?
\begin{enumerate}[label=\Alph*)]
\item 30
\item 60
\item 90
\item 120
\end{enumerate}

\item Piyoda A punktdan 30 km uzoqlashgach, piyodanikidan 5 marta katta tezlik bilan velosipedchi yo‘lga chiqdi. 
Agar ular B punktga bir paytda kirib borgan bo‘lsalar, punktlar orasidagi masofa qanday (km)?
\begin{enumerate}[label=\Alph*)]
\item 20
\item 22.4
\item 25
\item 37.5
\item 36
\end{enumerate}

\item A va B punktlardan bir-biriga qarab ikki velosipedchi yo‘lga chiqdi. 
Ular B punktga 30 km qolganda uchrashishadi. 
Manzilga borib qaytishda esa A punktdan 18 km uzoqlikda uchrashishdi. 
Punktlar orasidagi masofani toping (km).
\begin{enumerate}[label=\Alph*)]
\item 48
\item 60
\item 68
\item 72
\item 78
\end{enumerate}

\item Ikki jismning harakat tenglamalari $x_1 = 25 + t, \quad x_2 = 85 - 0.5t$ ko‘rinishga ega. Ularning uchrashish vaqti va joyi aniqlansin.
\begin{enumerate}[label=\Alph*)]
\item 40 s dan so‘ng, koordinata boshidan 65 m masofada
\item 18 s dan so‘ng, koordinata boshidan 114 m masofada
\item 10 s dan so‘ng, koordinata boshidan 5 m masofada
\item 30 s dan so‘ng, koordinata boshidan 60 m masofada
\item TJY
\end{enumerate}

\item Harakat tenglamalari $x_1 = 2 + 3t, \quad x_2 = 1 - 2t$ bo‘lgan moddiy nuqtalar vaqt boshidan necha sekund o‘tganda uchrashadilar?
\begin{enumerate}[label=\Alph*)]
\item ular uchrashmaydilar
\item 1
\item 2
\item 3
\item 4
\end{enumerate}

\item Yerdan 135 m balandlikdan uchib ketayotgan dronni razvedkachi durbin yordamida kuzatmoqda. 
Bir paytda dron vertikalga nisbatan $45^\circ$ burchak ostida ko‘rindi. 
Shundan 6 s o‘tgach $53^\circ$ burchak ostida ko‘rindi. 
Dronning tezligini toping (km/h). 
$\sin53^\circ=0.8$, $\cos53^\circ=0.6$.
\begin{enumerate}[label=\Alph*)]
\item 15
\item 34
\item 48
\item 27
\item 9
\end{enumerate}

\item Jism gorizont bilan $60^\circ$ burchak hosil qilgan qiyalikdan $20\ \text{m/s}$ tezlik bilan yuqoriga tortilmoqda. 
Jismning vertikal va gorizontal yo‘nalishlardagi tezliklari qanday (m/s)?
\begin{enumerate}[label=\Alph*)]
\item 17.3; 10
\item 10; 14.1
\item 10; 17.3
\item 14.1; 10
\item TJY
\end{enumerate}

\item Moddiy nuqta XOY tekisligida $x = 11 + 1.5t, \quad y = 3 + 2t$ tenglamalar bilan harakat qilmoqda. Tezligini aniqlang (m/s).
\begin{enumerate}[label=\Alph*)]
\item 4
\item 3.5
\item 5
\item 5.25
\item 2.5
\end{enumerate}

\item Moddiy nuqtaning harakati $x = 7 + 12t, \quad y = 5 + 16t$ tenglamalar bilan berilgan. Dastlabki 3 s ichida qancha ko‘chadi (m)?
\begin{enumerate}[label=\Alph*)]
\item 14
\item 19
\item 10
\item 13
\item 60
\end{enumerate}

\item Moddiy nuqta $ x = 2 + t, \quad y = 3 + 2t$ tenglamalar bilan harakat qiladi. Trayektoriya tenglamasini toping.
\begin{enumerate}[label=\Alph*)]
\item $y=2x$
\item $y=2x-1$
\item $y=2x+2$
\item $y=x+1$
\item $y=2x+3$
\end{enumerate}


\end{enumerate}



\begin{table}[h!]
\centering
\begin{tabular}{|c|c|c|c|c|c|c|c|c|c|}
\hline
1a & 2b & 3a & 4b & 5a & 6c & 7e & 8c & 9d & 10d \\ \hline
11a & 12a & 13d & 14a & 15e & 16e & 17b &  &  &  \\ \hline
\end{tabular}
\end{table}

